%&latex 
\documentclass[12pt]{article} 
\usepackage{mathrsfs} 
\usepackage{amsmath,amsfonts,amssymb}
\newcommand{\PR}{\mathscr{P}}
\newcommand{\DR}{\mathscr{D}}
%\usepackage{preambule}
\begin{document} 

\gdef\mpxshipout{\shipout\hbox\bgroup
  \setbox0=\hbox\bgroup}
\gdef\stopmpxshipout{\egroup  \dimen0=\ht0 \advance\dimen0\dp0
  \dimen1=\ht0 \dimen2=\dp0
  \setbox0=\hbox\bgroup
    \box0
    \ifnum\dimen0>0 \vrule width1sp height\dimen1 depth\dimen2 
    \else \vrule width1sp height1sp depth0sp\relax
    \fi\egroup
  \ht0=0pt \dp0=0pt \box0 \egroup}
\mpxshipout% line 83 ./polyn.tex
$1$%
\stopmpxshipout
\mpxshipout% line 84 ./polyn.tex
$i$%
\stopmpxshipout
\mpxshipout% line 85 ./polyn.tex
$-1$%
\stopmpxshipout
\mpxshipout% line 86 ./polyn.tex
$-i$%
\stopmpxshipout
\mpxshipout% line 87 ./polyn.tex
$ $%
\stopmpxshipout
\mpxshipout% line 88 ./polyn.tex
$0$%
\stopmpxshipout
\mpxshipout% line 98 ./polyn.tex
$e^{i{\pi \over 3}}$%
\stopmpxshipout
\mpxshipout% line 99 ./polyn.tex
$e^{i{2\pi \over 3}}$%
\stopmpxshipout
\mpxshipout% line 100 ./polyn.tex
$e^{i{\pi}}$%
\stopmpxshipout
\mpxshipout% line 101 ./polyn.tex
$e^{i{4\pi \over 3}}$%
\stopmpxshipout
\mpxshipout% line 102 ./polyn.tex
$e^{i{5\pi \over 3}}$%
\stopmpxshipout
\mpxshipout% line 103 ./polyn.tex
$e^{i{6\pi \over 3}}$%
\stopmpxshipout
\mpxshipout% line 221 ./polyn.tex
${\pi \over 4}$%
\stopmpxshipout
\mpxshipout% line 222 ./polyn.tex
${\pi \over 2}$%
\stopmpxshipout
\mpxshipout% line 223 ./polyn.tex
${3\pi \over 4}$%
\stopmpxshipout
\mpxshipout% line 224 ./polyn.tex
${\pi}$%
\stopmpxshipout
\mpxshipout% line 225 ./polyn.tex
$-{3\pi \over 4}$%
\stopmpxshipout
\mpxshipout% line 226 ./polyn.tex
$-{\pi \over 2}$%
\stopmpxshipout
\mpxshipout% line 227 ./polyn.tex
$-{\pi \over 4}$%
\stopmpxshipout
\mpxshipout% line 228 ./polyn.tex
${\sqrt{2} \over 2}$%
\stopmpxshipout
\mpxshipout% line 229 ./polyn.tex
$-{\sqrt{2} \over 2}$%
\stopmpxshipout
\mpxshipout% line 230 ./polyn.tex
${\sqrt{2} \over 2}$%
\stopmpxshipout
\mpxshipout% line 231 ./polyn.tex
$-{\sqrt{2} \over 2}$%
\stopmpxshipout
\mpxshipout% line 261 ./polyn.tex
$0$%
\stopmpxshipout
\mpxshipout% line 262 ./polyn.tex
$1 \over2$%
\stopmpxshipout
\mpxshipout% line 263 ./polyn.tex
$- {1 \over 2}$%
\stopmpxshipout
\mpxshipout% line 264 ./polyn.tex
$\pi$%
\stopmpxshipout
\mpxshipout% line 266 ./polyn.tex
{${\pi \over 3}$}%
\stopmpxshipout
\mpxshipout% line 267 ./polyn.tex
{${2\pi \over 3}$}%
\stopmpxshipout
\mpxshipout% line 268 ./polyn.tex
{$-{2\pi \over 3}$}%
\stopmpxshipout
\mpxshipout% line 269 ./polyn.tex
{$-{\pi \over 3}$}%
\stopmpxshipout
\mpxshipout% line 270 ./polyn.tex
${\sqrt{3} \over 2}$%
\stopmpxshipout
\mpxshipout% line 271 ./polyn.tex
$-{\sqrt{3} \over 2}$%
\stopmpxshipout
\mpxshipout% line 292 ./polyn.tex
${\pi \over 6}$%
\stopmpxshipout
\mpxshipout% line 293 ./polyn.tex
${\pi \over 2}$%
\stopmpxshipout
\mpxshipout% line 294 ./polyn.tex
${5\pi \over 6}$%
\stopmpxshipout
\mpxshipout% line 295 ./polyn.tex
$-{5\pi \over 6}$%
\stopmpxshipout
\mpxshipout% line 296 ./polyn.tex
$-{\pi \over 2}$%
\stopmpxshipout
\mpxshipout% line 297 ./polyn.tex
$-{\pi \over 6}$%
\stopmpxshipout
\mpxshipout% line 298 ./polyn.tex
${\sqrt{3} \over 2}$%
\stopmpxshipout
\mpxshipout% line 299 ./polyn.tex
$-{\sqrt{3} \over 2}$%
\stopmpxshipout
\mpxshipout% line 300 ./polyn.tex
${1 \over 2}$%
\stopmpxshipout
\mpxshipout% line 301 ./polyn.tex
$-{1 \over 2}$%
\stopmpxshipout
\mpxshipout% line 318 ./polyn.tex
${\pi \over 4}$%
\stopmpxshipout
\mpxshipout% line 319 ./polyn.tex
${\pi \over 2}$%
\stopmpxshipout
\mpxshipout% line 320 ./polyn.tex
${3\pi \over 4}$%
\stopmpxshipout
\mpxshipout% line 321 ./polyn.tex
${\pi}$%
\stopmpxshipout
\mpxshipout% line 322 ./polyn.tex
$-{3\pi \over 4}$%
\stopmpxshipout
\mpxshipout% line 323 ./polyn.tex
$-{\pi \over 2}$%
\stopmpxshipout
\mpxshipout% line 324 ./polyn.tex
$-{\pi \over 4}$%
\stopmpxshipout
\mpxshipout% line 325 ./polyn.tex
${\sqrt{2} \over 2}$%
\stopmpxshipout
\mpxshipout% line 326 ./polyn.tex
$-{\sqrt{2} \over 2}$%
\stopmpxshipout
\mpxshipout% line 327 ./polyn.tex
${\sqrt{2} \over 2}$%
\stopmpxshipout
\mpxshipout% line 328 ./polyn.tex
$-{\sqrt{2} \over 2}$%
\stopmpxshipout
\mpxshipout% line 383 ./polyn.tex
$1$%
\stopmpxshipout
\mpxshipout% line 384 ./polyn.tex
$i$%
\stopmpxshipout
\mpxshipout% line 385 ./polyn.tex
$-1$%
\stopmpxshipout
\mpxshipout% line 386 ./polyn.tex
$-i$%
\stopmpxshipout
\mpxshipout% line 387 ./polyn.tex
$ $%
\stopmpxshipout
\mpxshipout% line 388 ./polyn.tex
$0$%
\stopmpxshipout
\mpxshipout% line 398 ./polyn.tex
$e^{i{\pi \over 2}}$%
\stopmpxshipout
\mpxshipout% line 399 ./polyn.tex
$e^{i\pi}$%
\stopmpxshipout
\mpxshipout% line 400 ./polyn.tex
$e^{i{3\pi \over 2}}$%
\stopmpxshipout
\mpxshipout% line 401 ./polyn.tex
$e^{i{4\pi \over 2}}$%
\stopmpxshipout
\mpxshipout% line 488 ./polyn.tex
$1$%
\stopmpxshipout
\mpxshipout% line 489 ./polyn.tex
$i$%
\stopmpxshipout
\mpxshipout% line 490 ./polyn.tex
$-1$%
\stopmpxshipout
\mpxshipout% line 491 ./polyn.tex
$-i$%
\stopmpxshipout
\mpxshipout% line 492 ./polyn.tex
$ $%
\stopmpxshipout
\mpxshipout% line 493 ./polyn.tex
$0$%
\stopmpxshipout
\mpxshipout% line 503 ./polyn.tex
$e^{i{2\pi \over 3}}$%
\stopmpxshipout
\mpxshipout% line 504 ./polyn.tex
$e^{i{4\pi \over 3}$%
\stopmpxshipout
\mpxshipout% line 505 ./polyn.tex
$e^{i{6\pi \over 3}}$%
\stopmpxshipout
\end{document}
